\chapter{Time-Domain Analysis}
\label{chap:time}

\begin{learningobjective}
\begin{itemize}[leftmargin=*]
    \item Judge stability of linear systems using eigenvalue analysis, Lyapunov methods, and characteristic polynomials.
    \item Relate transient specifications (rise time, settling time, overshoot) to system poles and damping ratios.
    \item Follow complete numerical examples that compute performance metrics and validate them with Python simulations.
\end{itemize}
\end{learningobjective}

\section{Concept Map for Beginners}
\begin{conceptroadmap}
\begin{enumerate}[leftmargin=*]
    \item Begin with qualitative questions—does the system settle, oscillate, or diverge?—and tie them directly to eigenvalues and characteristic polynomials.
    \item Translate performance requirements into pole locations using a small library of formulas for damping ratio and natural frequency.
    \item Verify every analytical claim with the provided simulations, reinforcing algebra by observing the same effect in plots and logged metrics.
\end{enumerate}
\end{conceptroadmap}

\section{Concept--Math--Engineering Loop}
\begin{table}[h]
    \centering
    \small
    \renewcommand{\arraystretch}{1.2}
    \begin{tabularx}{\linewidth}{|>{\raggedright\arraybackslash}p{0.28\linewidth}|>{\raggedright\arraybackslash}p{0.34\linewidth}|>{\raggedright\arraybackslash}X|}
        \hline
        \textbf{Concept} & \textbf{Mathematics} & \textbf{Engineering Practice} \\
        \hline
        Stability assessment & Eigenvalues, Lyapunov equation, Routh--Hurwitz & Run eigenvalue checks and Lyapunov solves in Python to classify dynamics quickly. \\
        \hline
        Transient specifications & Second-order step response formulas, damping ratio $\zeta$ & Adjust $\zeta$ and $\omega_n$ in \texttt{tutorial\_time\_domain} to see predicted overshoot and settling time. \\
        \hline
        Controller interpretation & Root locus algebra, PID difference equations & Compare analytic root locus steps with the scripted plots and PID experiment logs. \\
        \hline
    \end{tabularx}
    \caption{How Module~\ref{chap:time} balances theory, computation, and hands-on validation.}
\end{table}

\section{Stability by Eigenvalues and Lyapunov Functions}
\begin{beginnerguide}
    \item[\textbf{What?}] Analytical tools for determining whether system trajectories decay, persist, or diverge.
    \item[\textbf{Why?}] Stability is the first safety check—controllers are useless if the closed loop explodes.
    \item[\textbf{When?}] Before tuning performance; confirm the open- or closed-loop poles sit in the stable region.
    \item[\textbf{How?}] Inspect eigenvalues of $A$ or solve Lyapunov equations to obtain energy-like certificates.
    \item[\textbf{Example}] Compute $P$ from $A^\top P + P A = -Q$ for a second-order system and verify numerically that $V(x)=x^\top P x$ decreases along trajectories.
\end{beginnerguide}
\subsection{Step-by-Step Derivation}
The homogeneous system $\dot{x} = A x$ has solution $x(t) = \mathrm{e}^{A t} x(0)$. Eigenvalues of $A$ govern behaviour:
\begin{itemize}[leftmargin=*]
    \item $\Re(\lambda_i) < 0$ for all $i$ $\Rightarrow$ asymptotically stable.
    \item Eigenvalues on the imaginary axis with no repetitions $\Rightarrow$ marginally stable.
    \item Any $\Re(\lambda_i) > 0$ or repeated imaginary-axis pole $\Rightarrow$ unstable.
\end{itemize}
For a constructive certificate, choose $Q = Q^\top \succ 0$ and solve the Lyapunov equation
\[
    A^\top P + P A = -Q.
\]
If $P \succ 0$, the quadratic function $V(x) = x^\top P x$ decreases along trajectories, proving asymptotic stability.

\subsection{Code Mapping}
\begin{table}[h]
    \centering
    \small
    \renewcommand{\arraystretch}{1.2}
    \begin{tabularx}{\linewidth}{|>{\raggedright\arraybackslash}p{4cm}|X|X|}
        \hline
        \textbf{Python} & \textbf{Formula} & \textbf{Explanation} \\
        \hline
        \texttt{np.linalg.eigvals(A)} & $\lambda_i(A)$ & Computes eigenvalues used to categorise stability. \\
        \hline
        \texttt{scipy.linalg.solve\_continuous\_lyapunov(A.T, -Q)} & $A^\top P + P A = -Q$ & Returns $P$; check positive definiteness with \texttt{np.all(np.linalg.eigvals(P)>0)}. \\
        \hline
    \end{tabularx}
    \caption{Eigenvalue and Lyapunov stability checks.}
\end{table}

\section{Characteristic Polynomials and Routh--Hurwitz}
\begin{beginnerguide}
    \item[\textbf{What?}] A tabular test for assessing polynomial stability without explicit root finding.
    \item[\textbf{Why?}] High-order characteristic equations arise frequently; Routh--Hurwitz provides a fast manual or symbolic check.
    \item[\textbf{When?}] When designing classical controllers that alter pole locations via parameter tuning.
    \item[\textbf{How?}] Construct the Routh table coefficients and ensure the first column remains positive.
    \item[\textbf{Example}] Apply the procedure to a cubic closed-loop polynomial and confirm the conditions using \texttt{sympy.routh\_array()}.
\end{beginnerguide}
\subsection{Derivation Workflow}
For a unity-feedback loop with open-loop transfer $L(s)$, the closed-loop characteristic polynomial is $D(s) + N(s)$. Routh--Hurwitz builds a tabular test without root finding. Example for $s^3 + a_1 s^2 + a_2 s + a_3$:
\[
\begin{array}{c|cc}
s^3 & 1 & a_2 \\
s^2 & a_1 & a_3 \\
s^1 & \dfrac{a_1 a_2 - a_3}{a_1} & 0 \\
s^0 & a_3 & 
\end{array}
\]
Every first-column entry must be positive. The template generalises to higher orders; replicate the calculation symbolically with SymPy to cement the algebra.

\subsection{Code Mapping}
\begin{table}[h]
    \centering
    \small
    \renewcommand{\arraystretch}{1.2}
    \begin{tabularx}{\linewidth}{|>{\raggedright\arraybackslash}p{4cm}|X|X|}
        \hline
        \textbf{Python} & \textbf{Formula} & \textbf{Explanation} \\
        \hline
        \texttt{sympy.routh\_array(poly)} & Routh table entries & Automates the eliminations; inspect intermediate determinants. \\
        \hline
        \texttt{numpy.poly(stable\_roots)} & $D(s)$ reconstruction & Useful for cross-checking with explicit pole placement. \\
        \hline
    \end{tabularx}
    \caption{Routh--Hurwitz computations in Python.}
\end{table}

\section{Second-Order Response Metrics}
\begin{beginnerguide}
    \item[\textbf{What?}] Closed-form relations between damping ratio, natural frequency, and time-domain specifications.
    \item[\textbf{Why?}] These approximations convert qualitative requirements (overshoot, settling time) into numeric design targets.
    \item[\textbf{When?}] Early in controller tuning to estimate suitable pole locations or PID gains.
    \item[\textbf{How?}] Use the standard second-order step response equations and solve for $\zeta$ and $\omega_n$ that meet specs.
    \item[\textbf{Example}] Given $M_p \le 10\%$ and $t_s \le 1.5$~s, solve for $\zeta \approx 0.59$, $\omega_n \ge 4.5$~rad/s, then simulate to validate.
\end{beginnerguide}
\subsection{Derivation}
Starting from
\[
    G(s) = \frac{\omega_n^2}{s^2 + 2 \zeta \omega_n s + \omega_n^2},
\]
use inverse Laplace transforms to obtain
\[
    y(t) = 1 - \frac{1}{\sqrt{1 - \zeta^2}} e^{-\zeta \omega_n t} \sin(\omega_d t + \phi),
\]
with $\omega_d = \omega_n \sqrt{1 - \zeta^2}$. Differentiating this step response yields classic metrics:
\begin{align*}
    t_r &\approx \frac{1.8}{\omega_n} && 0.4 \leq \zeta \leq 0.8,\\
    t_p &= \frac{\pi}{\omega_d},\\
    M_p &= e^{-\pi \zeta / \sqrt{1 - \zeta^2}},\\
    t_s &\approx \frac{4}{\zeta \omega_n} && \text{(2\% criterion)}.
\end{align*}

\subsection{Worked Specification Example}
Design for $t_s \leq 1.5$~s and $M_p \leq 10\%$:
\begin{enumerate}[leftmargin=*]
    \item Solve $0.10 = e^{-\pi \zeta/\sqrt{1-\zeta^2}} \Rightarrow \zeta \approx 0.59$.
    \item Impose $t_s \approx 4/(\zeta \omega_n) \leq 1.5$ to get $\omega_n \geq 4.5$~rad/s.
    \item Target dominant poles at $-\zeta \omega_n \pm \mathrm{j} \omega_d$ with $\omega_d = \omega_n \sqrt{1-\zeta^2} \approx 3.7$~rad/s.
\end{enumerate}

\subsection{Code Mapping}
\begin{table}[h]
    \centering
    \small
    \renewcommand{\arraystretch}{1.2}
    \begin{tabularx}{\linewidth}{|>{\raggedright\arraybackslash}p{4cm}|X|X|}
        \hline
        \textbf{Python} & \textbf{Formula} & \textbf{Explanation} \\
        \hline
        \texttt{control.step\_info(G)} & $t_r, t_p, M_p, t_s$ & Computes metrics from the transfer function. \\
        \hline
        \texttt{tutorial\_time\_domain} in \texttt{scripts/tutorials.py} & $y(t)$ vs. design targets & Visualises the effect of changing $\zeta$, $\omega_n$. \\
        \hline
    \end{tabularx}
    \caption{Second-order performance checks.}
\end{table}

\section{Root-Locus Construction}
\begin{beginnerguide}
    \item[\textbf{What?}] A graphical method showing how closed-loop poles move as a scalar gain varies.
    \item[\textbf{Why?}] It reveals stability margins and helps choose gains without exhaustive simulation.
    \item[\textbf{When?}] During classical controller design, especially when adjusting proportional gain or placing lead/lag networks.
    \item[\textbf{How?}] Plot the locus using pole/zero locations, angle criteria, and asymptote formulas.
    \item[\textbf{Example}] Sketch the locus for $K/(s(s^2+0.4s+4))$, then confirm with \texttt{control.root\_locus()} to see the poles leave the real axis.
\end{beginnerguide}
\subsection{Manual Derivation}
For $L(s) = K \frac{N(s)}{D(s)}$, root-locus points satisfy $1 + L(s) = 0$. The steepest steps mirror the derivation style in \texttt{kalman\_deduction.tex}:
\begin{enumerate}[leftmargin=*]
    \item Identify pole/zero structure of $L(s)$.
    \item Apply angle criterion $\angle L(s) = (2k+1)\pi$ to determine locus on the real axis.
    \item Compute asymptotes: $\sigma = (\sum \text{poles} - \sum \text{zeros})/(n - m)$, $\theta_k = \frac{(2k+1)\pi}{n-m}$.
    \item Evaluate $\frac{dK}{ds} = 0$ to locate breakaway points.
\end{enumerate}

\subsection{Example}
For $L(s) = K/(s(s^2 + 0.4 s + 4))$, differentiating $K(s) = -s(s^2 + 0.4s + 4)$ gives $3s^2 + 0.8 s + 4 = 0$, whose complex roots show that the locus leaves the real axis immediately.

\subsection{Code Mapping}
\begin{table}[h]
    \centering
    \small
    \renewcommand{\arraystretch}{1.2}
    \begin{tabularx}{\linewidth}{|>{\raggedright\arraybackslash}p{4cm}|X|X|}
        \hline
        \textbf{Python} & \textbf{Formula} & \textbf{Explanation} \\
        \hline
        \texttt{control.root\_locus(L)} & Locus plot & Visualises steps 1--4 graphically. \\
        \hline
        \texttt{tutorial\_time\_domain}\\\texttt{generate\_root\_locus()} & Annotated asymptotes & Recreates the analytical derivation numerically. \\
        \hline
    \end{tabularx}
    \caption{Root-locus derivation and verification.}
\end{table}

\section{PID Controllers: Time-Domain Interpretation}
\begin{beginnerguide}
    \item[\textbf{What?}] The time-domain view of proportional-integral-derivative control and its discrete implementation.
    \item[\textbf{Why?}] Understanding how each term shapes error evolution informs tuning and anti-windup strategies.
    \item[\textbf{When?}] Whenever you deploy PID loops or need a baseline controller before moving to advanced methods.
    \item[\textbf{How?}] Derive the difference equation, implement it with saturation guards, and tune gains using response metrics.
    \item[\textbf{Example}] Use \texttt{PIDController.print\_incremental\_coefficients()} to inspect discrete coefficients and simulate the mass--spring--damper case to see their effect.
\end{beginnerguide}
\subsection{Incremental Implementation}
The discrete incremental PID law is
\[
    u[k] = u[k-1] + K_p (e[k] - e[k-1]) + K_i T_s e[k] + \frac{K_d}{T_s} (e[k] - 2 e[k-1] + e[k-2]).
\]
Derivative filtering and anti-windup strategies (clamping, back-calculation) ensure robustness to quantization and saturation.

\subsection{Simulation Case Study}
\begin{lstlisting}[language=python]
from scripts.experiments import EXPERIMENTS

pid_result = EXPERIMENTS["pid"].runner()
print(pid_result.performance_metrics())
\end{lstlisting}
Typical metrics include a steady-state error of $\approx 6\times 10^{-3}$ and overshoot near $75\%$ for the aggressive tuning. Compare these with second-order predictions to tune $K_p,K_i,K_d$ systematically.

\subsection{Code Mapping}
\begin{table}[h]
    \centering
    \small
    \renewcommand{\arraystretch}{1.2}
    \begin{tabularx}{\linewidth}{|>{\raggedright\arraybackslash}p{4cm}|X|X|}
        \hline
        \textbf{Python} & \textbf{Formula} & \textbf{Explanation} \\
        \hline
        \texttt{PIDController.step()} & Incremental PID law & Implements the difference equation and anti-windup hooks. \\
        \hline
        \texttt{tutorial\_time\_domain} & $y(t)$, $u(t)$ plots & Shows how gain changes alter $t_s$, $M_p$, and actuator effort. \\
        \hline
    \end{tabularx}
    \caption{PID controller equations versus code.}
\end{table}

\section{Disturbance and Noise Rejection}
\begin{beginnerguide}
    \item[\textbf{What?}] Analysis of how feedback loops attenuate load disturbances and amplify measurement noise.
    \item[\textbf{Why?}] Controllers must balance robustness against disturbances with immunity to sensor imperfections.
    \item[\textbf{When?}] After satisfying nominal performance, evaluate disturbance/noise behaviour before hardware deployment.
    \item[\textbf{How?}] Examine sensitivity functions $S$ and complementary sensitivity $T$, augment simulations with injected disturbances.
    \item[\textbf{Example}] Modify \texttt{tutorial\_time\_domain} to add a step disturbance at $t=2$~s and compare the resulting $y(t)$ across PID tunings.
\end{beginnerguide}
\subsection{Derivation}
For unity feedback,
\[
    Y(s) = \frac{P(s)}{1 + P(s) C(s)} D(s) + \frac{P(s) C(s)}{1 + P(s) C(s)} R(s).
\]
Low-frequency disturbances are attenuated when $|PC| \gg 1$; sensor noise amplification grows with high-frequency loop gain, highlighting the need for derivative filtering.

\subsection{Case Study}
Modify \texttt{tutorial\_time\_domain} to inject a step disturbance at $t=2$~s and Gaussian sensor noise with variance $10^{-4}$. Compare the resulting plots under aggressive and conservative PID tunings to visualise the sensitivity functions $S = 1/(1+PC)$ and $T = PC/(1+PC)$.

\section{Workflow Summary and Practice}
\begin{beginnerguide}
    \item[\textbf{What?}] A checklist for applying time-domain techniques end-to-end.
    \item[\textbf{Why?}] Summaries help ensure you validate assumptions, tune gains responsibly, and document results reproducibly.
    \item[\textbf{When?}] After completing the chapter and whenever you start a new design iteration.
    \item[\textbf{How?}] Follow the enumerated steps: map specs to poles, verify stability, cross-check with simulation, and log results.
    \item[\textbf{Example}] Run \tutorialcmd[--output-dir figures]{time}, annotate rise time/overshoot from the plots, and compare with second-order predictions.
\end{beginnerguide}
\begin{enumerate}[leftmargin=*]
    \item Translate specifications into pole targets using the second-order formulas and root-locus insights.
    \item Verify algebraic stability tests (Lyapunov, Routh--Hurwitz) with Python to avoid arithmetic mistakes.
    \item Validate against simulations using \tutorialcmd[--output-dir figures]{time}; inspect \texttt{figures/time\_pid\_responses.png} and \texttt{figures/time\_pid\_controls.png}.
    \item Log innovations or residuals in experiments to tune disturbance/noise models before moving to frequency-domain design.
\end{enumerate}
Refer back to Appendix~\ref{chap:tooling} for the reproducible plotting commands and \texttt{make} targets that automate these checks.
\IfFileExists{../figures/time_pid_responses.png}{
    \begin{figure}[h]
        \centering
        \includegraphics[width=0.75\linewidth]{../figures/time_pid_responses.png}
        \caption{PID step responses for multiple tunings generated via \tutorialcmd{time}.}
        \label{fig:pid-responses}
    \end{figure}
}{
    \begin{figure}[h]
        \centering
        \fbox{\parbox{0.7\linewidth}{Run \tutorialcmd{time} to create the multi-case PID response plot, then recompile.}}
        \caption{Placeholder for the multi-case PID step-response plot.}
        \label{fig:pid-responses}
    \end{figure}
}
\IfFileExists{../figures/time_pid_controls.png}{
    \begin{figure}[h]
        \centering
        \includegraphics[width=0.75\linewidth]{../figures/time_pid_controls.png}
        \caption{Control effort comparison across PID tunings.}
        \label{fig:pid-controls}
    \end{figure}
}{
    \begin{figure}[h]
        \centering
        \fbox{\parbox{0.7\linewidth}{Generate the control-effort plot with \tutorialcmd{time} to visualise actuator demand.}}
        \caption{Placeholder for the PID control-effort plot.}
        \label{fig:pid-controls}
    \end{figure}
}

\begin{takeaway}
\begin{itemize}[leftmargin=*]
    \item Stability analysis combines algebraic tests (Routh--Hurwitz) with geometric intuition (root locus, eigenvalues).
    \item Performance formulas for second-order systems provide quick design rules which must be cross-checked against simulations.
    \item Python-backed experiments translate equations into numerical intuition, highlighting effects of nonidealities.
\end{itemize}
\end{takeaway}

\section{Module Review}
\begin{itemize}[leftmargin=*]
    \item \textbf{Stability}: eigenvalue locations, Lyapunov equation $A^\top P + P A = -Q$.
    \item \textbf{Routh--Hurwitz}: all first-column entries positive $\Leftrightarrow$ closed-loop stability.
    \item \textbf{Second-order metrics}: $t_s \approx 4/(\zeta \omega_n)$, $M_p = e^{-\pi \zeta / \sqrt{1-\zeta^2}}$.
    \item \textbf{Root locus}: asymptote centroid $\sigma$ and angles $\theta_k$ describe pole migration.
    \item \textbf{PID implementation}: incremental difference equation with anti-windup safeguards.
    \item \textbf{Disturbance rejection}: $S = 1/(1 + PC)$, $T = PC/(1 + PC)$ quantify load disturbance versus noise amplification.
\end{itemize}

\section{Keyword Review}
\begin{vocabulary}
\begin{itemize}[leftmargin=*]
    \item \textbf{Asymptotic stability}: all eigenvalues with negative real parts (continuous time) or magnitude less than one (discrete time).
    \item \textbf{Lyapunov function}: energy-like scalar used to certify stability; decreases along trajectories.
    \item \textbf{Characteristic polynomial}: determinant $\det(sI - A)$ (or closed-loop denominator) whose roots are system poles.
    \item \textbf{Damping ratio $\zeta$}: dimensionless measure controlling overshoot and oscillation in second-order systems.
    \item \textbf{Root locus}: plot showing closed-loop pole trajectories as a gain varies.
    \item \textbf{Sensitivity $S$}: transfer function mapping disturbances to output, highlighting robustness limits.
\end{itemize}
\end{vocabulary}
